% Part: model-theory
% Chapter: models-of-arithmetic
% Section: standard-models

\documentclass[../../../include/open-logic-section]{subfiles}

\begin{document}

\olfileid{mod}{mar}{stm}
\olsection{د حساب معياري مدلونه}

د حساب ژبه~$\Lang{L_A}$ په څرګند ډول د شمېرو، په ځانګړي ډول د
طبيعي شمېرو، په هکله ده۔ نو «معياري مدل»~$\Struct{N}$ ځانګړے دے:
همدا هغه مدل دے چې موږ ئې په هکله خبرې کول غواړو۔ خو په منطق کښې
زموږ پام ډېر ځله يوازې جوړښتي ځانګړتياوو ته وي، او هر دوه آيزومورف
!!{structure}s دغه ځانګړتياوې ګډې لري۔ نو لږ پراخ فکر کولے شو او
هر هغه !!{structure} چې له~$\Struct{N}$ سره آيزومورف وي «معياري»
وبولو۔

\begin{defn}
د~$\Lang{L_A}$ يو !!{structure} \emph{معياري} دے، که له~$\Struct{N}$
سره آيزومورف وي۔
\end{defn}

\begin{prop}
\ollabel{prop:standard-domain} که !!a{structure}~$\Struct{M}$ معياري
وي، نو تعريف ساحه ئې د معياري عددنښو د ارزښتونو سټ دے، يعنې:
\[
\Domain{M} = \Setabs{\Value{\num{n}}{M}}{n \in \Nat}
\]
\end{prop}

\begin{proof}
څرګنده ده چې هر $\Value{\num{n}}{M} \in \Domain{M}$ دے۔ يوازې دا
ښودل پاتې دي چې هر $x \in \Domain{M}$ د کوم~$n$ دپاره
$\Value{\num{n}}{M}$ دے۔ څرنګه چې $\Struct{M}$ معياري دے، له
$\Struct{N}$ سره آيزومورف دے۔ فرض کړئ $g\colon \Nat \to \Domain{M}$
آيزومورفيزم وي۔ بيا $g(n) = g(\Value{\num{n}}{N}) =
\Value{\num{n}}{M}$۔ خو د هر $x \in \Domain{M}$ دپاره داسې~$n \in
\Nat$ شته چې $g(n) = x$ وي، ځکه $g$ !!{surjective} ده۔
\end{proof}

\begin{explain}
که د~$\Lang{L_A}$ يو جوړښت~$\Struct{M}$ معياري وي، د~!!{domain}
ټول غړي ئې د معياري عددنښو $\num{0}$، $\num{1}$، $\num{2}$،
\dots په نوم يادېدے شي، يعنې د $\Obj{0}$، $\Obj{0}'$،
$\Obj{0}''$ او داسې نورو ترمونو په وسيله۔ البته دا معنا نۀ لري چې د
$\Domain{M}$ !!{element}s پخپله شمېرې دي؛ يوازې دومره ده چې د هغو
په نښه کولے شو، لکه څنګه چې په~$\Domain{N}$ کښې شمېرې نښه کوو۔
\end{explain}

\begin{prob}
وښيئ چې د~\olref[mod][mar][stm]{prop:standard-domain} معکوس ناسم دے؛
يعنې داسې يو !!a{structure}~$\Struct{M}$ ورکړئ چې
$\Domain{M} = \Setabs{\Value{\num{n}}{M}}{n \in \Nat}$ وي، خو له
$\Struct{N}$ سره آيزومورف نۀ وي۔
\end{prob}

\begin{prop}
\ollabel{prop:thq-standard}
که $\Sat{M}{\Th{Q}}$ او $\Domain{M} =
\Setabs{\Value{\num{n}}{M}}{n \in \Nat}$ وي، نو $\Struct{M}$ معياري دے۔
\end{prop}

\begin{proof}
بايد وښيو چې $\Struct{M}$ له~$\Struct{N}$ سره آيزومورف دے۔ تابع
$g\colon \Nat \to \Domain{M}$ داسې ونيسئ: $g(n) =
\Value{\num{n}}{M}$۔ د فرض له مخې $g$ !!{surjective} ده۔ دا
!!{injective} هم ده: هر کله چې $n \neq m$ وي، $\Th{Q} \Proves
\eq/[\num{n}][\num{m}]$۔ نو له $\Sat{M}{\Th{Q}}$ څخه
$\Sat{M}{\eq/[\num{n}][\num{m}]}$، هر کله چې $n \neq m$۔ په پايله
کښې $n \neq m$ وي، نو $\Value{\num{n}}{M} \neq
\Value{\num{m}}{M}$، يعنې $g(n) \neq g(m)$۔

اوس بايد دا هم وڅېړو چې $g$ آيزومورفيزم دے۔
\begin{enumerate}
\item $g(\Assign{\Obj{0}}{N}) = g(0)$، ځکه $\Assign{\Obj{0}}{N} = 0$۔
  د~$g$ له تعريفه $g(0) = \Value{\num{0}}{M}$۔ خو $\num{0}$ هماغه
  $\Obj{0}$ دے، او د داسې ترم ارزښت چې !!a{constant} وي، هغه څه دے
  چې !!{structure} ئې دغه !!{constant} ته ټاکي، يعنې $\Value{\Obj{0}}{M} =
  \Assign{\Obj{0}}{M}$۔ نو $g(\Assign{\Obj{0}}{N}) =
  \Assign{\Obj{0}}{M}$، لکه اړين چې و۔
\item $g(\Assign{\prime}{N}(n)) = g(n+1)$، ځکه په~$\Struct{N}$ کښې
  $\prime$ پر~$\Nat$ د جانشين تابع ده۔ بيا د~$g$ له تعريفه
  $g(n+1) = \Value{\num{n+1}}{M}$۔ خو $\num{n+1}$ له~$\num{n}'$
  سره يو ترم دے، نو $\Value{\num{n+1}}{M} =
  \Value{\num{n}'}{M}$۔ د ارزښت د تابعې له تعريفه دا
  $= \Assign{\prime}{M}(\Value{\num{n}}{M})$ دے۔ څرنګه چې
  $\Value{\num{n}}{M} = g(n)$، نو
  $g(\Assign{\prime}{N}(n)) = \Assign{\prime}{M}(g(n))$۔
\item $g(\Assign{+}{N}(n,m)) = g(n+m)$، ځکه په~$\Struct{N}$ کښې $+$
  د~$\Nat$ د جمع تابع ده۔ بيا د~$g$ له تعريفه
$g(n+m) = \Value{\num{n+m}}{M}$۔
  خو $\Th{Q} \Proves \num{n+m} = (\num{n} + \num{m})$، نو
  $\Value{\num{n+m}}{M} = \Value{\num{n}+\num{m}}{M}$۔ د ارزښت د
  تابعې له تعريفه دا
  $= \Assign{+}{M}(\Value{\num{n}}{M},\Value{\num{m}}{M})$ دے۔
  څرنګه چې $\Value{\num{n}}{M} = g(n)$ او
  $\Value{\num{m}}{M} = g(m)$، نو
  $g(\Assign{+}{N}(n, m)) = \Assign{+}{M}(g(n), g(m))$۔
\item $g(\Assign{\times}{N}(n, m)) = \Assign{\times}{M}(g(n), g(m))$:
  تمرين۔
\item $\tuple{n,m} \in \Assign{<}{N}$ هله او يوازې هله چې $n < m$۔
  که $n < m$ وي، نو $\Th{Q} \Proves \num{n} < \num{m}$ او
  $\Sat{M}{\num{n} < \num{m}}$ هم لرو۔ نو
  $\tuple{\Value{\num{n}}{M}, \Value{\num{m}}{M}} \in
  \Assign{<}{M}$، يعنې $\tuple{g(n), g(m)} \in \Assign{<}{M}$۔
  که $n \not< m$ وي، نو $\Th{Q} \Proves \lnot \num{n} < \num{m}$ او
  په پايله کښې $\Sat/{M}{\num{n} < \num{m}}$۔ نو، لکه مخکې،
  $\tuple{g(n), g(m)} \notin \Assign{<}{M}$۔ په ګډه:
  $\tuple{n,m} \in \Assign{<}{N}$ هله او يوازې هله چې
  $\tuple{g(n), g(m)} \in \Assign{<}{M}$۔
\end{enumerate}
\end{proof}

\begin{explain}
تابع~$g$ د~$\Lang{L_A}$ د هر بل !!{structure}~$\Struct{M}$ د تعريف
ساحې ته له~$\Nat$ څخه د نقشې تر ټولو څرګنده لار ده، ځکه هر داسې
$\Struct{M}$ هغه !!{element}s لري چې $\num{0}$، $\num{1}$،
$\num{2}$ او داسې نورو په نومونو يادېږي۔ نو د حيرانتيا خبره نۀ ده چې
که $\Struct{M}$ د~$\num{n}$ په باب لږ تر لږه ځينې بنسټيزې ويناوې
هماغسې رښتونې کړي لکه~$\Struct{N}$ او $g$ هم دوه اړخيزه يو په يو
وي، نو $g$ آيزومورفيزم ګرځي۔ په حقيقت کښې که $\Domain{M}$ له هغو
!!{element}s پرته چې $\num{n}$ ئې نوموي بل څه نۀ لري، نو همدا
يوازينے آيزومورفيزم دے۔
\end{explain}

\begin{prop}
\ollabel{prop:thq-unique-iso} که $\Struct{M}$ معياري وي، نو د
\olref{prop:thq-standard} له ثبوت څخه $g$ له~$\Struct{N}$ څخه
$\Struct{M}$ ته يوازينے آيزومورفيزم دے۔
\end{prop}

\begin{proof}
فرض کړئ $h\colon \Nat \to \Domain{M}$ له~$\Struct{N}$ څخه
$\Struct{M}$ ته آيزومورفيزم وي۔ د~$n$ پر بنسټ استقرا ښيو چې $g=h$۔
که $n=0$ وي، نو د~$g$ له تعريفه $g(0)=\Assign{\Obj{0}}{M}$۔ خو
څرنګه چې $h$ آيزومورفيزم دے، $h(0)=h(\Assign{\Obj{0}}{N}) =
\Assign{\Obj{0}}{M}$، نو $g(0)=h(0)$۔

اوس د $n+1$ حالت راواخلئ:
\begin{align*}
  g(n+1) & = \Value{\num{n+1}}{M} \text{ د $g$ له تعريفه}\\
  & = \Value{\num{n}'}{M} \text{ ځکه $\num{n+1}\ident \num{n}'$}\\
  & = \Assign{\prime}{M}(\Value{\num{n}}{M})
    \text{ د $\Value{t'}{M}$ له تعريفه}\\
  & = \Assign{\prime}{M}(g(n)) \text{ د $g$ له تعريفه}\\
  & = \Assign{\prime}{M}(h(n)) \text{ د استقرايي فرض له مخې}\\
  & = h(\Assign{\prime}{N}(n)) \text{ ځکه $h$ آيزومورفيزم دے}\\
  & = h(n+1)
\end{align*}
\end{proof}

\begin{explain}
د هر !!{denumerable} سټ~$M$ دپاره د~$\Nat$ او~$M$ تر منځ يوه
!!a{bijection} شته، نو هر داسې~$M$ په بالقوه ډول د معياري
مدل~$\Struct{M}$ !!{domain} کېدے شي۔ په حقيقت کښې، کله چې يو څيز
$z \in M$ او مناسبه تابع~$s$ د~$\Assign{\Obj{0}}{M}$ او
$\Assign{\prime}{M}$ په توګه وټاکئ، د $+$، $\times$ او $<$ تفسيرونه
لا له مخکې ټاکل شوي وي۔ يوازې هغه تابعې~$s\colon M \to M \setminus
\{z\}$ په معياري مدل کښې د~$\Assign{\prime}{M}$ په توګه مناسبې دي چې
!!{injective} او !!{surjective} دواړه وي۔ د~$s$ د قيمتونو ساحه
$z$ لرلے نۀ شي، ځکه که داسې واے نو
$\lforall[x][\eq/[\Obj 0][x']]$ به دروغ واے۔ دا !!{sentence}
په~$\Struct{N}$ کښې رښتونې ده، نو $\Struct{M}$ ئې هم بايد رښتونې کړي۔ تابع~$s$ بايد
!!{injective} وي، ځکه په~$\Struct{N}$ کښې د جانشين
تابع~$\Assign{\prime}{N}$ همداسې ده، او د هغې
$\Assign{\prime}{N}$ !!{injective} توب په~$\Struct{N}$ کښې د
!!a{sentence} له خوا بيانېږي۔ دا بايد
!!{surjective} هم وي، ځکه که نه نو داسې~$x \in M \setminus \{z\}$
به وي چې د~$s$ د قيمتونو په ساحه کښې نۀ وي، يعنې دا !!{sentence}
$\lforall[x][(\eq[x][\Obj 0] \lor
\lexists[y][\eq[y'][x]])]$ به په~$\Struct{M}$ کښې دروغ وي، حال دا چې
په~$\Struct{N}$ کښې رښتونې ده۔
\textbf{د ژباړې د نسخې سمون (OLMOD-005):} سرچينې په وروستۍ جمله کښې
د تابعې د «قيمتونو ساحې» پر ځاے د «تعريف ساحه» ليکلې وه، حال دا
چې تابع s پر ټول M تعريف شوې ده۔ دلته د جملې له فورمول او د
پرښتياوالي له شرط سره سم د قيمتونو ساحه ليکل شوې ده۔
\end{explain}

\end{document}
